\begin{proof}[Proof of \cref{obj:lem:original-score-covariance-ledger}]
\begin{enumerate}
\item Suppose first that \(n<4\). Since \(n\ge2\), the zero-score convention in \cref{obj:def:explicit-score-ledger,obj:def:calibration-handle} gives
\[
Z_{b,0}=Z_{b,A}=0,
\qquad
\Lambda_{n,v}=0.
\]
Both scalar coefficients are square integrable, and their variances and cross-covariance vanish. This proves all the asserted bounds in this case.
% lean: CausalSmith.Stat.FinitepHomogeneityDensegamma.covarianceLedger_small_sample

\item Henceforth \(n\ge4\). Use the ranks, cutoffs, and ledgers of \cref{obj:def:explicit-score-ledger}, and the blocks of \cref{obj:def:blocks}. In particular,
\[
s=s_n=\lfloor n/2\rfloor\ge2,
\qquad
M>0,\quad K>0,\quad M\mid K,
\qquad
T_0\ge1,
\]
and, in the multiresolution branch,
\[
R_j=2^jM,\qquad M\mid R_{j-1},\qquad T_j\ge1
\quad(1\le j\le L).
\]
The divisibility follows because the selected ranks are powers of two and \(K\ge M\).

Fix \(u\in\mathbb R^M\). Write
\[
\mathcal O=[0,1]\times\{0,1\}\times\mathbb R,
\qquad
O_i=(X_i,A_i,Y_i)\in\mathcal O,
\]
and let \(O,O'\) be independent records with law \(P\). For \(x\in[0,1]\), define the scalar histogram feature by
\[
f_u(x)=u^TF_M(x).
\]
The disjoint histogram cells in \cref{obj:def:projection}, together with the uniform covariate law, give
\[
\int_0^1 f_u(x)^2\,dx
=
\sum_{j=1}^{M}u_j^2
=
\|u\|_2^2,
\qquad
\mathbb E[f_u(X)^2]=\|u\|_2^2.
\]

Use a coefficient selector \(r\in\{0,1\}\), with
\[
Z_{b,r}^{\mathrm{coef}}
=
\begin{cases}
Z_{b,0},&r=0,\\
Z_{b,A},&r=1.
\end{cases}
\]
For every real cutoff \(T\ge1\), define
\[
\mathfrak m_{r,T}(o)
=
\begin{cases}
\ell_T(Y(o)),&r=0,\\
A(o),&r=1,
\end{cases}
\qquad
V_T=32T^{2-p},
\qquad o\in\mathcal O.
\]
Thus \(V_{T_0}=V_0\) and \(V_{T_j}=V_j\). For a pair kernel \(G\) from the score construction, define
\[
\mathscr E_{r,u}[G,T](o,o')
=
f_u(X(o))A(o)G(X(o),X(o'))\mathfrak m_{r,T}(o'),
\]
\[
\mathscr P_{r,u}[G,T](o,o')
=
\frac{
\mathscr E_{r,u}[G,T](o,o')
+
\mathscr E_{r,u}[G,T](o',o)
}{2},
\qquad o,o'\in\mathcal O.
\]
The selected scalar single-record and pair summands are
\[
h_{r,u}(o)=f_u(X(o))A(o)\mathfrak m_{r,T_0}(o),
\]
\[
g_{r,u}(o,o')
=
\begin{cases}
-\mathscr P_{r,u}[\Pi_K,T_0](o,o'),
&\text{single-correction branch},\\[1mm]
-\mathscr P_{r,u}[\Pi_M,T_0](o,o')
-\displaystyle\sum_{j=1}^{L}\mathscr P_{r,u}[D_j,T_j](o,o'),
&\text{multiresolution branch}.
\end{cases}
\]
Here \(A(o)^2=A(o)\), so the \(r=1\) single-record summand agrees with the treatment coefficient. Package these summands into
\[
k_{r,u}(o,o')
=
\frac{h_{r,u}(o)+h_{r,u}(o')}{2}
+
g_{r,u}(o,o').
\]
All these functions are measurable and bounded for the fixed finite tuning, hence square integrable.

For either block, symmetrization preserves the ordered-pair sum, and each index occurs \(s-1\) times in each position. Consequently, the coefficient definitions in \cref{obj:def:calibration-handle} and the score formulas in \cref{obj:def:score-family} give
\[
u^TZ_{b,r}^{\mathrm{coef}}
=
\frac1s\sum_{i\in\mathcal I_b}h_{r,u}(O_i)
+
\frac1{s(s-1)}
\sum_{\substack{i,j\in\mathcal I_b\\i\ne j}}
g_{r,u}(O_i,O_j)
=
\frac1{s(s-1)}
\sum_{\substack{i,j\in\mathcal I_b\\i\ne j}}
k_{r,u}(O_i,O_j).
\]
Empty increment sums have value zero throughout.
% lean: CausalSmith.Stat.FinitepHomogeneityDensegamma.ledgerScalarKernel_average

\item We first establish the variance identity for this pair-average representation. For any measurable symmetric square-integrable kernel
\(\mathscr H:\mathcal O^2\to\mathbb R\), define
\[
\overline{\mathscr H}=\mathbb E[\mathscr H(O,O')],
\qquad
\mathscr H^{(1)}(o)
=
\mathbb E[\mathscr H(o,O')]-\overline{\mathscr H},
\]
\[
\mathscr H^{(2)}(o,o')
=
\mathscr H(o,o')
-\overline{\mathscr H}
-\mathscr H^{(1)}(o)
-\mathscr H^{(1)}(o'),
\]
and
\[
\mathscr A_b(\mathscr H)
=
\frac1{s(s-1)}
\sum_{\substack{i,j\in\mathcal I_b\\i\ne j}}
\mathscr H(O_i,O_j).
\]
Conditional integration gives
\[
\mathbb E[\mathscr H^{(1)}(O)]=0,
\qquad
\mathbb E[\mathscr H^{(2)}(O,O')\mid O]
=
\mathbb E[\mathscr H^{(2)}(O,O')\mid O']
=
0.
\]
Expanding the square of
\(\mathscr H=\overline{\mathscr H}+\mathscr H^{(1)}+\mathscr H^{(1)}+\mathscr H^{(2)}\)
and integrating therefore yields
\[
\mathbb E[\mathscr H(O,O')^2]
=
\overline{\mathscr H}^{\,2}
+
2\mathbb E[\mathscr H^{(1)}(O)^2]
+
\mathbb E[\mathscr H^{(2)}(O,O')^2].
\]
Indeed, the two singleton terms have zero product expectation by independence, and their products with the canonical term have zero expectation by conditional integration.

Counting the occurrences of each singleton term gives
\[
\mathscr A_b(\mathscr H)-\overline{\mathscr H}
=
\frac2s\sum_{i\in\mathcal I_b}\mathscr H^{(1)}(O_i)
+
\mathscr A_b(\mathscr H^{(2)}).
\]
The singleton sum has variance
\[
\operatorname{Var}\!\left(
\frac2s\sum_{i\in\mathcal I_b}\mathscr H^{(1)}(O_i)
\right)
=
\frac4s\operatorname{Var}(\mathscr H^{(1)}(O)).
\]
For distinct indices within each pair, conditional centering and independence imply
\[
\operatorname{Cov}\!\left(
\mathscr H^{(2)}(O_i,O_j),
\mathscr H^{(2)}(O_{i'},O_{j'})
\right)
=
\begin{cases}
\mathbb E[\mathscr H^{(2)}(O,O')^2],
&(i',j')=(i,j)\text{ or }(j,i),\\
0,&\text{otherwise}.
\end{cases}
\]
The zero case includes pairs sharing exactly one index: integrate first in a coordinate that occurs in only one pair. There are \(s(s-1)\) ordered pairs and two contributing ordered pairs for each of them. Thus
\[
\operatorname{Var}(\mathscr A_b(\mathscr H^{(2)}))
=
\frac{2}{s(s-1)}
\mathbb E[\mathscr H^{(2)}(O,O')^2].
\]
The same conditional-centering argument gives
\[
\operatorname{Cov}\!\left(
\frac2s\sum_{i\in\mathcal I_b}\mathscr H^{(1)}(O_i),
\mathscr A_b(\mathscr H^{(2)})
\right)=0.
\]
Each evaluation block has marginal law \(P^{\otimes s}\) under \(P^{\otimes n}\), so
\[
\operatorname{Var}_{P^{\otimes n}}(\mathscr A_b(\mathscr H))
=
\frac4s\operatorname{Var}_P(\mathscr H^{(1)})
+
\frac{2}{s(s-1)}
\mathbb E[\mathscr H^{(2)}(O,O')^2].
\]
% lean: CausalSmith.Stat.FinitepHomogeneityDensegamma.evalBlockPairAverage_variance

\item We next bound the singleton channel. Admissibility gives \(1<p\le2\). For \(T\ge1\), clipping satisfies the pointwise bounds
\[
|y-\ell_T(y)|\le |y|^pT^{1-p},
\qquad
\ell_T(y)^2\le |y|^pT^{2-p}.
\]
Integrating under the conditional moment restriction in \cref{obj:ass:raw-moment} gives, for each treatment arm and almost every covariate value,
\[
\left|
\int\{y-\ell_T(y)\}\,Q_{a,P}(dy\mid x)
\right|
\le10T^{1-p},
\qquad
\int\ell_T(y)^2\,Q_{a,P}(dy\mid x)
\le10T^{2-p}.
\]
Define the conditional mark mean, for \(x\in[0,1]\), by
\[
\overline{\mathfrak m}_{r,T}(x)
=
e_P(x)\int\mathfrak m_{r,T}(x,1,y)\,Q_{1,P}(dy\mid x)
+
\{1-e_P(x)\}
\int\mathfrak m_{r,T}(x,0,y)\,Q_{0,P}(dy\mid x).
\]
The preceding bounds and \(0\le e_P\le1\) imply
\[
\mathbb E[\mathfrak m_{0,T}(O)^2\mid X=x]
\le10T^{2-p}\le V_T.
\]
For the treatment mark,
\[
\mathbb E[\mathfrak m_{1,T}(O)^2\mid X=x]
=e_P(x)\le1\le V_T,
\qquad
\overline{\mathfrak m}_{1,T}(x)=e_P(x).
\]
In particular, for both selectors,
\[
\mathbb E[\mathfrak m_{r,T}(O)^2\mid X=x]\le V_T,
\qquad
|\overline{\mathfrak m}_{r,T}(x)|\le\sqrt{V_T}
\quad\text{almost everywhere},
\]
where the second inequality is conditional Cauchy--Schwarz.

Since \(|A|\le1\), the single-record summand obeys
\[
h_{r,u}(O)^2
\le
2f_u(X)^2\mathfrak m_{r,T_0}(O)^2.
\]
Conditioning on \(X\) and using the feature isometry yields
\[
\mathbb E[h_{r,u}(O)^2]
\le2V_0\|u\|_2^2
\le B_h^2\|u\|_2^2,
\qquad
B_h=2\sqrt{V_0}.
\]
% lean: CausalSmith.Stat.FinitepHomogeneityDensegamma.multiresScalarSingle_energy_le

\item For any selected projection or difference kernel \(G\), and any integrable function \(\varphi:[0,1]\to\mathbb R\), use the operator notation
\[
(G\varphi)(x)=\int_0^1G(x,x')\varphi(x')\,dx'.
\]
Nested histogram cells imply
\[
G(x,x')f_u(x')=G(x,x')f_u(x)
\]
for \(G=\Pi_K,\Pi_M\), and for every \(G=D_j\). Integrating the two symmetrized legs separately therefore gives
\[
\mathbb E[\mathscr P_{r,u}[G,T](o,O')]
=
\frac12f_u(X(o))
\left[
A(o)(G\overline{\mathfrak m}_{r,T})(X(o))
+
\mathfrak m_{r,T}(o)(Ge_P)(X(o))
\right].
\]
For a positive histogram rank \(J_{\mathrm{aux}}\) divisible by \(M\), cell averaging gives
\[
|\Pi_{J_{\mathrm{aux}}}\overline{\mathfrak m}_{r,T}|
\le\sqrt{V_T},
\qquad
|\Pi_{J_{\mathrm{aux}}}e_P|\le1.
\]
Consequently,
\[
\left|
\mathbb E[\mathscr P_{r,u}[\Pi_{J_{\mathrm{aux}}},T](o,O')]
\right|
\le
\frac12|f_u(X(o))|
\left(\sqrt{V_T}+|\mathfrak m_{r,T}(o)|\right).
\]
Using
\[
\frac14\left(\sqrt{V_T}+|\mathfrak m_{r,T}(o)|\right)^2
\le
\frac12\left(V_T+\mathfrak m_{r,T}(o)^2\right)
\]
and then conditioning on \(X\), we obtain
\[
\mathbb E\!\left[
\left\{
\mathbb E[\mathscr P_{r,u}[\Pi_{J_{\mathrm{aux}}},T](O,O')\mid O]
\right\}^{\!2}
\right]
\le V_T\|u\|_2^2.
\]

The augmented kernel has singleton projection
\[
k_{r,u}^{(1)}(o)
=
\frac{h_{r,u}(o)-\mathbb E[h_{r,u}(O)]}{2}
+
g_{r,u}^{(1)}(o).
\]
Centering decreases second moment, so the triangle inequality in \(L_2(P)\) gives
\[
2\|k_{r,u}^{(1)}\|_{L_2(P)}
\le
\sqrt{\mathbb E[h_{r,u}(O)^2]}
+
2\sqrt{
\mathbb E\!\left[
\left\{\mathbb E[g_{r,u}(O,O')\mid O]\right\}^{2}
\right]}.
\]
Here
\[
\operatorname{Var}_P(k_{r,u}^{(1)})
=
\|k_{r,u}^{(1)}\|_{L_2(P)}^2
\]
because its mean is zero.

In the single-correction branch, the projection estimate with
\(J_{\mathrm{aux}}=K\) and \(T=T_0\) gives
\[
\mathbb E\!\left[
\left\{\mathbb E[g_{r,u}(O,O')\mid O]\right\}^{2}
\right]
\le B_{g,\mathrm{single}}^2\|u\|_2^2,
\qquad
B_{g,\mathrm{single}}=\sqrt{V_0}.
\]
It follows that
\[
4\operatorname{Var}_P(k_{r,u}^{(1)})
\le
(B_h+2B_{g,\mathrm{single}})^2\|u\|_2^2
=
(4\sqrt{V_0})^2\|u\|_2^2
\le
(16\sqrt{V_0})^2\|u\|_2^2
=
L_1^2\|u\|_2^2.
\]
% lean: CausalSmith.Stat.FinitepHomogeneityDensegamma.ledgerScalarKernel_singleton_variance_le

\item In the multiresolution branch, define
\[
s_0=\min\{\alpha,\beta,\gamma\}\in(0,1],
\qquad
f_{\mathrm{mean}}(x)=m_{0,P}(x)+e_P(x)\tau_P(x),
\qquad x\in[0,1].
\]
The smoothness and boundedness restrictions in
\cref{obj:ass:propensity-smoothness,obj:ass:baseline-smoothness,obj:ass:effect-smoothness,obj:ass:overlap,obj:ass:baseline-cap,obj:ass:effect-cap}
give
\[
|f_{\mathrm{mean}}(x)|\le2
\]
and
\[
\begin{aligned}
|f_{\mathrm{mean}}(x)-f_{\mathrm{mean}}(x')|
&\le
|m_{0,P}(x)-m_{0,P}(x')|
+
|e_P(x)|\,|\tau_P(x)-\tau_P(x')|\\
&\quad+
|e_P(x)-e_P(x')|\,|\tau_P(x')|\\
&\le
20|x-x'|^\beta
+
15|x-x'|^\gamma
+
20|x-x'|^\alpha\\
&\le55|x-x'|^{s_0}.
\end{aligned}
\]
For \(x=x'\) both sides vanish; for \(0<|x-x'|\le1\), the last inequality follows because each exponent is at least \(s_0\).

To apply the histogram approximation bound with constant \(20\), define
\[
\widetilde f_{\mathrm{mean}}(x)=\frac{20}{55}f_{\mathrm{mean}}(x).
\]
Then
\[
\|\widetilde f_{\mathrm{mean}}\|_\infty\le\frac{40}{55}\le20,
\qquad
|\widetilde f_{\mathrm{mean}}(x)-\widetilde f_{\mathrm{mean}}(x')|
\le20|x-x'|^{s_0}.
\]
Averaging over the histogram cell containing \(x\), whose diameter is at most \(J_{\mathrm{aux}}^{-1}\), gives, for every positive integer \(J_{\mathrm{aux}}\),
\[
|\Pi_{J_{\mathrm{aux}}}\widetilde f_{\mathrm{mean}}(x)
-\widetilde f_{\mathrm{mean}}(x)|
\le20J_{\mathrm{aux}}^{-s_0}.
\]
Thus, for \(1\le j\le L\),
\[
\begin{aligned}
|D_jf_{\mathrm{mean}}(x)|
&=
\frac{55}{20}
\left|
\Pi_{R_j}\widetilde f_{\mathrm{mean}}(x)
-
\Pi_{R_{j-1}}\widetilde f_{\mathrm{mean}}(x)
\right|\\
&\le
\frac{55}{20}
\left(20R_j^{-s_0}+20R_{j-1}^{-s_0}\right)
\le110R_{j-1}^{-s_0}.
\end{aligned}
\]

For \(T\ge1\), define on all of \([0,1]\)
\[
\mathfrak t_T(x)
=
f_{\mathrm{mean}}(x)-\overline{\mathfrak m}_{0,T}(x).
\]
The conditional mean versions in \cref{obj:def:model} imply, almost everywhere,
\[
\begin{aligned}
\mathfrak t_T(x)
&=
e_P(x)\int\{y-\ell_T(y)\}\,Q_{1,P}(dy\mid x)\\
&\quad+
\{1-e_P(x)\}
\int\{y-\ell_T(y)\}\,Q_{0,P}(dy\mid x).
\end{aligned}
\]
The clipping-tail estimate therefore gives
\[
\|\mathfrak t_T\|_{L_\infty(\lambda)}
\le10T^{1-p},
\qquad
|D_j\mathfrak t_{T_j}(x)|
\le20T_j^{1-p}.
\]
For the propensity, the same cell-averaging argument gives
\[
|\Pi_{J_{\mathrm{aux}}}e_P(x)-e_P(x)|
\le20J_{\mathrm{aux}}^{-\alpha},
\]
hence
\[
|D_je_P(x)|
\le20R_j^{-\alpha}+20R_{j-1}^{-\alpha}
\le40R_{j-1}^{-\alpha}=a_j.
\]
Combining the smoothness and tail bounds yields
\[
|D_j\overline{\mathfrak m}_{0,T_j}(x)|
\le110R_{j-1}^{-s_0}+20T_j^{1-p}=d_j.
\]

Substitution into the conditional pair formula gives the outcome-coefficient envelope
\[
\left|
\mathbb E[\mathscr P_{0,u}[D_j,T_j](o,O')]
\right|
\le
\frac12|f_u(X(o))|
\left(d_j+a_j|\mathfrak m_{0,T_j}(o)|\right).
\]
Since
\[
\frac14(d_j+a_j|\mathfrak m_{0,T_j}|)^2
\le
\frac12(d_j^2+a_j^2\mathfrak m_{0,T_j}^2),
\]
conditioning on \(X\) gives
\[
\mathbb E\!\left[
\left\{
\mathbb E[\mathscr P_{0,u}[D_j,T_j](O,O')\mid O]
\right\}^{\!2}
\right]
\le
\frac12(d_j^2+a_j^2V_j)\|u\|_2^2.
\]
For the treatment coefficient,
\(\overline{\mathfrak m}_{1,T_j}=e_P\) and
\(|\mathfrak m_{1,T_j}|\le1\), so
\[
\left|
\mathbb E[\mathscr P_{1,u}[D_j,T_j](o,O')]
\right|
\le a_j|f_u(X(o))|.
\]
Applying the same square-envelope inequality with constant term \(2a_j\) and zero mark term yields
\[
\mathbb E\!\left[
\left\{
\mathbb E[\mathscr P_{1,u}[D_j,T_j](O,O')\mid O]
\right\}^{\!2}
\right]
\le2a_j^2\|u\|_2^2.
\]
All \(a_j,d_j\) are nonnegative and \(V_j\ge1\). The two coefficient estimates fit the common bound because
\[
\frac12(d_j^2+a_j^2V_j)
\le\{2(d_j+a_j\sqrt{V_j})\}^2,
\qquad
2a_j^2
\le\{2(d_j+a_j\sqrt{V_j})\}^2.
\]
Define the nonnegative increment sum and conditional-pair budget by
\[
\mathcal S_{\mathrm{inc}}
=
\sum_{j=1}^{L}(d_j+a_j\sqrt{V_j}),
\qquad
B_{g,\mathrm{multi}}
=
\sqrt{V_0}+2\mathcal S_{\mathrm{inc}}.
\]
The initial projection estimate at rank \(M\), followed by the triangle inequality for the finite increment sum in \(L_2(P)\), gives
\[
\begin{aligned}
\left\|
\mathbb E[g_{r,u}(O,O')\mid O]
\right\|_{L_2(P)}
&\le
\sqrt{V_0}\|u\|_2
+
\sum_{j=1}^{L}2(d_j+a_j\sqrt{V_j})\|u\|_2\\
&=
B_{g,\mathrm{multi}}\|u\|_2.
\end{aligned}
\]
This includes \(L=0\), when the sum is zero. The singleton estimate for the augmented kernel now gives
\[
\begin{aligned}
4\operatorname{Var}_P(k_{r,u}^{(1)})
&\le
(B_h+2B_{g,\mathrm{multi}})^2\|u\|_2^2\\
&=
(4\sqrt{V_0}+4\mathcal S_{\mathrm{inc}})^2\|u\|_2^2\\
&\le
(16\sqrt{V_0}+8\mathcal S_{\mathrm{inc}})^2\|u\|_2^2
=
L_1^2\|u\|_2^2.
\end{aligned}
\]
% lean: CausalSmith.Stat.FinitepHomogeneityDensegamma.ledgerScalarKernel_singleton_variance_le

\item We next bound the canonical channel, again treating the two branches separately. The definitions of the projections give
\[
\overline{k_{r,u}}
=
\mathbb E[h_{r,u}(O)]+\overline{g_{r,u}},
\]
and cancellation of the single-record terms gives
\[
k_{r,u}^{(2)}(o,o')=g_{r,u}^{(2)}(o,o')
\quad\text{almost everywhere}.
\]
The pair-energy identity therefore implies
\[
\begin{aligned}
\mathbb E[k_{r,u}^{(2)}(O,O')^2]
&=
\mathbb E[g_{r,u}^{(2)}(O,O')^2]\\
&\le
\overline{g_{r,u}}^{\,2}
+
2\mathbb E[g_{r,u}^{(1)}(O)^2]
+
\mathbb E[g_{r,u}^{(2)}(O,O')^2]\\
&=
\mathbb E[g_{r,u}(O,O')^2].
\end{aligned}
\]

To bound this raw pair energy, the histogram kernels have exact row identities
\[
\int_0^1\Pi_{J_{\mathrm{aux}}}(x,x')^2\,dx'
=
J_{\mathrm{aux}}
\quad(J_{\mathrm{aux}}\in\mathbb N,\ J_{\mathrm{aux}}>0),
\]
\[
\int_0^1
\Pi_{R_j}(x,x')\Pi_{R_{j-1}}(x,x')\,dx'
=
R_{j-1},
\]
and hence
\[
\int_0^1D_j(x,x')^2\,dx'
=
R_j+R_{j-1}-2R_{j-1}
=
R_{j-1}
\quad(1\le j\le L).
\]
The first identity follows because a projection row equals its rank on one cell of reciprocal length. For the cross-product identity, the finer cell lies inside the coarser cell, so integrating their product gives the coarser rank.

Conditioning on the covariate of \(O'\), the mark bound gives, for every selected \(G,T\),
\[
\mathbb E[G(x,X(O'))^2\mathfrak m_{r,T}(O')^2]
\le
V_T\int_0^1G(x,x')^2\,dx'.
\]
Thus, integrating first in \(O'\) and using \(A(O)^2\le1\),
\[
\mathbb E[\mathscr E_{r,u}[\Pi_{J_{\mathrm{aux}}},T](O,O')^2]
\le
J_{\mathrm{aux}}V_T\,\mathbb E[f_u(X)^2]
=
J_{\mathrm{aux}}V_T\|u\|_2^2,
\]
\[
\mathbb E[\mathscr E_{r,u}[D_j,T_j](O,O')^2]
\le
R_{j-1}V_j\|u\|_2^2.
\]
Symmetrization satisfies
\[
\mathscr P_{r,u}[G,T](o,o')^2
\le
\frac{
\mathscr E_{r,u}[G,T](o,o')^2
+
\mathscr E_{r,u}[G,T](o',o)^2
}{2}.
\]
The two legs have equal integrated energy under \(P\otimes P\). Consequently,
\[
\mathbb E[\mathscr P_{r,u}[\Pi_{J_{\mathrm{aux}}},T](O,O')^2]
\le
J_{\mathrm{aux}}V_T\|u\|_2^2,
\]
\[
\mathbb E[\mathscr P_{r,u}[D_j,T_j](O,O')^2]
\le
R_{j-1}V_j\|u\|_2^2
\le
R_jV_j\|u\|_2^2.
\]

In the single-correction branch, this gives
\[
\begin{aligned}
\mathbb E[k_{r,u}^{(2)}(O,O')^2]
&\le
\mathbb E[g_{r,u}(O,O')^2]\\
&\le
KV_0\|u\|_2^2\\
&\le
(4\sqrt{KV_0})^2\|u\|_2^2
=
L_2^2\|u\|_2^2.
\end{aligned}
\]
In the multiresolution branch, the triangle inequality in \(L_2(P\otimes P)\) gives
\[
\|g_{r,u}\|_{L_2(P\otimes P)}
\le
\left(
\sqrt{MV_0}
+
\sum_{j=1}^{L}\sqrt{R_jV_j}
\right)\|u\|_2.
\]
Every summand is nonnegative, including the empty-sum case. Therefore
\[
\begin{aligned}
\mathbb E[k_{r,u}^{(2)}(O,O')^2]
&\le
\left(
\sqrt{MV_0}+\sum_{j=1}^{L}\sqrt{R_jV_j}
\right)^2\|u\|_2^2\\
&\le
\left[
4\left(
\sqrt{MV_0}+\sum_{j=1}^{L}\sqrt{R_jV_j}
\right)
\right]^2\|u\|_2^2\\
&=
L_2^2\|u\|_2^2.
\end{aligned}
\]
% lean: CausalSmith.Stat.FinitepHomogeneityDensegamma.ledgerScalarKernel_canonical_energy_le

\item Apply the pair-average variance identity to \(k_{r,u}\). Since \(s\ge2\), both sample factors are positive, and the two projection estimates imply
\[
\begin{aligned}
\operatorname{Var}_{P^{\otimes n}}(u^TZ_{b,r}^{\mathrm{coef}})
&=
\frac4s\operatorname{Var}_P(k_{r,u}^{(1)})
+
\frac{2}{s(s-1)}
\mathbb E[k_{r,u}^{(2)}(O,O')^2]\\
&\le
\frac{L_1^2}{s}\|u\|_2^2
+
\frac{2L_2^2}{s(s-1)}\|u\|_2^2\\
&=
\Lambda_{n,v}\|u\|_2^2.
\end{aligned}
\]
The finite-sum representation also gives
\[
u^TZ_{b,r}^{\mathrm{coef}}\in L_2(P^{\otimes n}).
\]
Taking \(r=0\) and \(r=1\) proves square integrability and the two variance bounds.
% lean: CausalSmith.Stat.FinitepHomogeneityDensegamma.original_score_covariance_ledger

\item Finally, the covariance ledger is nonnegative:
\[
\Lambda_{n,v}
=
\frac{L_1^2}{s}
+
\frac{2L_2^2}{s(s-1)}
\ge0
\qquad(n\ge4),
\]
and its value is zero for \(n=2,3\). For arbitrary \(u,z\in\mathbb R^M\), Cauchy--Schwarz for the centered square-integrable coefficients gives
\[
\begin{aligned}
\left|
\operatorname{Cov}_{P^{\otimes n}}
(u^TZ_{b,0},z^TZ_{b,A})
\right|
&\le
\sqrt{\operatorname{Var}_{P^{\otimes n}}(u^TZ_{b,0})}\,
\sqrt{\operatorname{Var}_{P^{\otimes n}}(z^TZ_{b,A})}\\
&\le
\sqrt{\Lambda_{n,v}\|u\|_2^2}\,
\sqrt{\Lambda_{n,v}\|z\|_2^2}\\
&=
(\sqrt{\Lambda_{n,v}})^2\|u\|_2\|z\|_2\\
&=
\Lambda_{n,v}\|u\|_2\|z\|_2.
\end{aligned}
\]
The square-root simplification applies also when either direction or the ledger is zero.
% lean: CausalSmith.Stat.FinitepHomogeneityDensegamma.ledger_cross_covariance_of_variance
\end{enumerate}
\end{proof}
